\section{A new axiom interface: an independent witness consequent}
\label{sec:witness}

The main logical experiment in this investigation changes an actual axiom
family, rather than merely the association of its emitter.  It also explains
how two individually unsafe-looking simplifications can be made compatible.
All the equivalences in this section concern finite numerical derivations;
they do not presuppose the consistency or the soundness of PA.

Let
\[
 p(a,b)=\frac{(a+b)(a+b+1)}2+b,\quad
 I(a,b)=4p(a,b),\quad M(a,b)=I(a,b)+1,
\]
\[
 U(d,B)=I(B,d)+2,\qquad E(a,b)=I(a,b)+3.
\]
The distinguished contradiction is the numerical code $f=1$.  Write
$A\to B=I(A,B)$ and $\neg A=I(A,1)$.  In particular, zero is a real
constructor alias: $I(0,0)=0$.  It is not legitimate to treat all these
constructors as disjoint free syntax after numerical interpretation.

The inherited calculus has modus ponens, theorem generalization
$B\mapsto U(d,B)$, the dominated implicational family
\[
 \rho_8(D,A,H,u)=((H\to u)\to D)\to((D\to H)\to(A\to H)),
 \qquad u\leq D,
\]
and the six heads
\begin{align*}
 O(D,H)&=\bigl(U(D,U(D,H)\to1)\to1\bigr)\to H,\\
 h_2(D,A,H)&=U(D,U(A,H))\to U(A,U(D,H)),\\
 h_3(D,A,H)&=M(H,D)\to(E(H,A)\to M(A,D)),\\
 h_4(D,A,H)&=E(H,D)\to(M(A,H)\to M(A,D)),\\
 h_5(D,H)&=U(D,E(H,D)\to1)\to1,\\
 h_6^Q(D,A,H)&=U(D,A\to H)\to(U(D,A)\to U(D,H)).
\end{align*}
It also retains the diagonal schema $\Gamma(B)=B\to U(B,B)$ and the
guarded pointed-separation head of the baseline.  Those two families,
their guards, and all inference rules are unchanged in the first variant.
The inherited restricted-$\rho_8$ development supplies uniform derived
$K$, $S$, and Peirce schemas before any appeal to $O$ or explosion.

Put $W(D,H)=U(D,E(H,D)\to1)$.  Replace the old fifth head $W(D,H)\to1$
by
\begin{equation}
             h_5^X(D,A,H)=W(D,H)\to A. \label{eq:h5X}
\end{equation}
Here $D,A,H$ range independently over all natural numbers.  The fact that
$A$ is independent of the equality label $H$ is essential.  The previously
investigated replacement $W(D,H)\to H$ does not provide the same
specialization argument for arbitrary $D,H$.

\begin{theorem}[Witness-consequent equivalence]
The old calculus and the calculus with \eqref{eq:h5X} in place of its fifth
head derive exactly the same numerical conclusions.  There are
primitive-recursive proof translations in both directions whose
correctness is provable in iPRA.
\end{theorem}

\begin{proof}
First recall why the old first head supplies explosion without assuming
that numerical code $1$ is already an available falsity rule.  For any
$A$, put $Y=U(0,U(0,A)\to1)$.  An instance of $K$ gives
$1\to(Y\to1)$, while $O(0,A)$ is $(Y\to1)\to A$.  Composition gives
$1\to A$.  Thus $W(D,H)\to1$ and $1\to A$ compose to give
\eqref{eq:h5X}.  In the reverse direction, specialize the independent
parameter to $A=1$.  The result is literally the old fifth-head code,
with no renaming of $D,H$ and no change to a quantifier binder.

Replace each changed axiom leaf by this fixed derivation and relocate
the indices of subsequent modus-ponens and generalization steps.  All
other leaves are copied.  Formula construction, finite concatenation,
index relocation and case distinction on a proof tag are primitive
recursive.  Correctness is quantifier-free induction on the Boolean
predicate saying that the finite output prefix checks and has the
specified last conclusion.  The output proof is a syntactic witness;
no reflection principle occurs.
\end{proof}

The delivered independent checker verifies a 13-step old-to-new template
and the one-step reverse specialization over the explicitly retained
derived $K,S$ basis.  The formulas are schematic terms with arbitrary
natural-number parameters.  Substituting their numerical meanings
preserves every checked modus-ponens equality, even when two distinct
schematic terms denote the same number.  This is stronger than testing
a finite range of parameter values, and weaker than a new proof-assistant
kernel verification of the entire inherited calculus.

\subsection{Explosion can move out of the first head}

The new witness family has a second useful property.  It supplies
explosion using only retained diagonal and quantifier-distribution
axioms, so the first head need not carry that job.

\begin{lemma}[Witness bootstrap]
The schemas $\Gamma,h_6^Q,h_5^X$, together with $K,S$, modus ponens and
closed generalization, uniformly derive $1\to A$.  Neither the old first
head nor Peirce's law is used.
\end{lemma}

\begin{proof}
Let $e=E(0,1)$ and $B=e\to1$.  The following are the significant
boundaries of the finite proof:
\begin{align*}
 &1\to U(1,1) &&\text{by }\Gamma(1),\\
 &1\to(e\to1) &&\text{by }K,\\
 &U(1,1\to B) &&\text{by theorem generalization},\\
 &U(1,1)\to U(1,B) &&\text{by }h_6^Q(1,1,B),\\
 &U(1,B)\to A &&\text{by }h_5^X(1,A,0).
\end{align*}
Compose the first, fourth and fifth lines.  The result is $1\to A$.
The only generalized formula is the unconditional $K$ theorem in the
second line; no generalization under an undischarged assumption occurs.
\end{proof}

The serialized certificate has 16 steps: one $\Gamma$ instance, one
generalization, one $h_6^Q$ instance, one $h_5^X$ instance, three $K$
instances, two $S$ instances and seven modus-ponens steps.  An independent
review replayed every line and checked that no $O$, old $h_5$, Peirce,
open assumption or propositional-tautology oracle is admitted.

\subsection{A coupled reduction of the quantifier head}

Set $X=U(D,U(D,H)\to1)$.  The old first head is
$O=(X\to1)\to H$.  Consider the join-shaped head
\begin{equation}
                   J(D,H)=(X\to H)\to H. \label{eq:J}
\end{equation}
Replacing $O$ by $J$ in isolation would not justify importing the old
calibration: the bootstrap of explosion has changed.  Replacing $O,h_5$
\emph{together} by $J,h_5^X$ does preserve the calculus.

\begin{theorem}[Coupled presentation]
The original numerical calculus and the calculus obtained by the two
replacements $O\rightsquigarrow J$ and $h_5\rightsquigarrow h_5^X$ are
uniformly interderivable in iPRA, with all other families and rules fixed.
\end{theorem}

\begin{proof}
For the new-to-old direction, the preceding lemma supplies $1\to H$.
Under the assumption $X\to1$, composition gives $X\to H$, and $J$ gives
$H$.  Discharge the assumption by the positive $K,S$ deduction
transformation.  This is exactly $O$.

For the old-to-new direction, assume $p:X\to H$.  Under the further
assumption $n:H\to1$, composition produces $X\to1$; $O$ then produces
$H$.  Discharge $n$ to obtain $(H\to1)\to H$ and apply Peirce's law
with $(P,Q)=(H,1)$.  Discharging $p$ gives $J$.  This direction does
not require a separate explosion lemma.  The old and new fifth heads
are translated as above.  Both transformations again replace leaves
by fixed templates and use Boolean induction on the coded proof list.
\end{proof}

The independently replayed $J$ certificates contain 95 steps for
$O\vdash J$ and 72 for the new calculus's derivation of $O$.  The latter
uses no old $O$, old $h_5$ or Peirce.  Their nonminimal sizes are
certificate sizes, not estimates of the size of a Turing-machine
implementation.

\subsection{Removing the remaining false constant from the first head}

There is a further conservative generalization:
\begin{equation}
 J^A(D,A,H)=\bigl(U(D,U(D,H)\to A)\to H\bigr)\to H.
 \label{eq:JA}
\end{equation}
Again $A$ is independent of $D,H$.  Specializing $A=1$ gives $J$ exactly.
Conversely put $X_1=U(D,U(D,H)\to1)$ and
$X_A=U(D,U(D,H)\to A)$.  Old explosion gives $1\to A$.
Positive composition, theorem generalization at $D$ and $h_6^Q$ give
$X_1\to X_A$.  Under the assumption $X_A\to H$, compose to get
$X_1\to H$ and apply $J$ to get $H$.  Discharging that assumption
gives \eqref{eq:JA}.  Thus replacing $O,h_5$ by $J^A,h_5^X$ again
preserves every numerical theorem, by primitive-recursive proof
translations verified in iPRA.  Its independently checked templates
have 358 steps from the old calculus to $J^A$ and 72 steps back to $O$.

This extra generalization does not invoke vacuity of $A$ at $D$.
It uses a generalized \emph{theorem} about implication and the retained
quantifier-distribution schema.  That distinction prevents a tempting
but invalid further shortcut: $U(D,E(H,D)\to A)\to A$ is not sound
for arbitrary $A$ that may contain the quantified variable.  For example,
take $D=0$, $H=1$, and $A=M(0,2)$ in the arithmetic bit interpretation,
with $x_0=1,x_1=0,x_2=1$.  Its antecedent holds because the only relevant
equality witness is zero and zero is a bit member of one, whereas its
consequent is false because one is not a bit member of itself.

There is a second coupled presentation using the already known
contraposed head $(H\to1)\to X$.  Its recovery of $O$ now has a valid
explosion premise, supplied by the 16-step bootstrap rather than by
adding another axiom head.  The delivered certificates have 204 and 131
steps for the two directions.  The contraposed head itself is prior
work; the witness bootstrap that makes the coupled replacement
legitimate is the new part of this construction.

\subsection{Consequences for source compilation}

In the AX11GH command, the independent parameters are
$A=c$, $D=d+x$ and $H=h+y$.  A selector invoked on a normalized store
has $d=h=0$, so all three parameters of \eqref{eq:h5X} can be selected
independently.  The old witness is selected by $c=1$, leaving $x=D$
and $y=H$ unrestricted.  Changing neighboring generated heads during
that selection is harmless: the selector preserves the desired head
and removes the others using its certified stack operations.

Both $h_5^X$ and $J$ have outer constructor $I$, so their codes are
divisible by four and cannot equal the halting code $1$.  The number
and order of emitted stack entries can therefore be retained.  In the
unpermuted form, replacing the fifth head's final
\texttt{push; incT; cons} by \texttt{copyA; cons} changes the AX word
from 190 to 189 macro tokens.  Replacing the corresponding false
consequent in $O$ by \texttt{copyH} changes the coupled word to 188.
Replacing the remaining inner false constant in that head by
\texttt{copyA} gives the $J^A,h_5^X$ word with 187 tokens.
These identities are checked in a free algebra against the complete
emitter, not just against the changed local token interval.

A smaller macro word is not a smaller Turing machine.  In particular,
removing a literal false constant can destroy a shared binary subtree
that formerly served several heads.  The physical compiler prices
shared reader controls and compatible history folds, and its entire
primitive control-flow graph still has to match the intended source
program.  The construction and search chapters state the exact
achieved state counts separately from these logical theorems.

\subsection{A family of further safe consequents}

The same proof gives a useful search principle.  Let $F$ be a fixed
primitive-recursive numerical expression, independent of $D,H$, and
suppose a fixed value $a_0$ has a retained derivation of
$F(a_0)\to1$.  Replacing the witness by $W(D,H)\to F(A)$ is
conservative when the old first head is retained: old explosion gives
every new instance, and specializing $A=a_0$ and composing with the
fixed refutation gives the old witness.

For $F(A)=A+1$, take $a_0=0$ and recover the old witness literally.
For $F(A)=A+k$ with $2\leq k\leq8$, choose $a_0=8-k$.
Here $8=I(0,1)$, and the retained theorem $0=\rho_8(0,0,0,0)$ supplies
$8\to1$ by positive implicational deduction.  Likewise
$F(A)=I(0,A)$ admits $a_0=1$, and $F(A)=M(A,A)$ admits $a_0=0$.
These provide justified emitter experiments, not automatic state-count
improvements.  The separate question of removing $O$ at the same time
requires an explosion proof for the chosen $F$; it must not be inferred
solely from this conservative-extension argument.
